\subsection{记号}

设 H 是 E 的一个超平面。回忆 E-H 有两个连通分支，称为由 H 界定的开半空间。它们的闭包称为由 H 界定的闭半空间。设 \(x, y \in E\)。如果 x 和 y 包含在由 H 界定的同一个开半空间内，或者等价地说，端点为 x 和 y 的闭线段不与 H 相交，则称 x 和 y 严格位于 H 的同侧；如果 x 属于由 H 界定的一个开半空间而 y 属于另一个，则称 x 和 y 位于 H 的相反两侧。若 \(x \in E\) 且 \(t \in T\)，如果对于所有 \(h + t\)，x 和 \(h \in H\) 严格位于 H 的同侧，则称 x 和 t 严格位于 H 的同侧。

设 A 是 E 的非空连通子集。对于 E 中不与 A 相交的任意超平面 H，用 \(D_{H}(A)\) 表示包含 A 的、由 H 界定的唯一开半空间。若 N 是 E 的一个超平面集，且其中没有一个与 A 相交，置

\[
D _ {\mathfrak {N}} (A) = \bigcap_ {H \in \mathfrak {N}} D _ {H} (A).\tag{1}
\]

如果 \(\mathrm{A}\) 由单点 \(a\) 组成，则我们用 \(\mathrm{D_H}(a)\) 和 \(\mathrm{D}_{\mathfrak{N}}(a)\) 代替 \(\mathrm{D_H}(\{a\})\) 和 \(\mathrm{D}_{\mathfrak{N}}(\{a\})\)。

